\section{Week 9}

\textbf{Important theorems} (that will come later):

\begin{theorem}[Cauchy's Theorem]
   If \( G  \) is a finite group and \(  p \) is a prime dividing the order of \( G  \), then \( G  \) has a element of order \( p  \).
\end{theorem}

\begin{theorem}[Sylow Theorem]
    If \( G  \) is a finite group of order \( p^{\alpha } \cdot m  \) where \( p  \) is a prime and \( p   \) does not divide \( m  \), then \( G  \) has a subgroup of order \( p^{\alpha} \) where \( \alpha \in \N \).
\end{theorem}

In the following, we will relax one of the assumptions found in Cauchy's theorem and assume that it is abelian.

\begin{prop}[21 Dummit and Foote]
   If \( G  \) is a finite abelian group and \( p  \) is a prime dividing \( | G |   \), then \( G  \) contains an element of order \( p  \).
\end{prop}

\begin{proof}
We will use strong induction to prove this result. (we will assume the result holds for all  groups with order is less than \( | G  |  \) and show this implies the result for \( | G  |  \)). Since \( | G  |  > 1  \), there exists an element \( x \in G  \) such that \( x \neq {1}_{G} \). If \( | G  |  = p  \), then we are done by Lagrange's theorem; that is, any group with order prime \( p  \) contains an element of order \( p  \). Assume that \( | G  |  > p  \). Suppose that \( p \mid | x  |  \). Hence,  
\[  | x  |  = p \cdot n \ \ \text{for some} \ n \in \Z.  \]
We know that \( | x^{n}  |  = \frac{ | x  |  }{  (| x  | , n )  } = \frac{ p \cdot n   }{ n  }   = p  \). Again, we have an element of order \( p  \) and we are done. We now assume that \( p \not\mid | x  |  \). Let \( \langle  x  \rangle = N  \). Since \( G  \) is abelian, this \( N  \) is normal in \( G  \). Hence, Lagrange's theorem implies that, we can say that 
\[  | G : N  |  = \Big|   G/ N  \Big| = \frac{ | G  |  }{  | N  |  }.   \]
Since \( N \neq {1}_{G}  \), we have that 
\[  \frac{ | G  |  }{  | N  |  }  < |  G  |. \]
Since \( p \not\mid | N  |  \), we have that \( p \mid |  G /  N  |  \). By our inductive assumption, we can conclude that there exists a \( y N \in G / N  \) such that 
\[  | y N  |  = p. \]
Since \( y N \neq 1 N  \), we have \( y \notin N  \) (in order for \( y \in N  \) we would need to raise to \( p  \) in order for it to be the identity). But \( y^{p} \in N  \) since \( (yN)^{p} = y^{p} N  = {1}_{G} N  \). Since \( \langle  y^{p} \rangle \leq N    \), we have that \( \langle  y \rangle \neq \langle  y^{p} \rangle \). That is, \( \langle   y^{p} \rangle < \langle y  \rangle \). Furthermore, \( | y  |  > | y^{p} |  \). We know \( | y^{p} |  = \frac{ | y  |   }{  (| y | , p ) }  \). But since \( p \mid | y  |  \), we have   
\[  | y^{p} |  = \frac{ | y  |  }{ p  }.  \]
Thus, \( p \mid | y  |  \) and we are done by our first case. Indeed, suppose for contradiction that \( (| y  | , p ) = 1  \), then from the above \( | y^{p} |  = | y  |  \) but this contradicts the assumption that \( | y^{p} |  < | y  |  \). This completes the induction and every finite abelian group with \( p \mid |  G  |  \) has an element of order \(  p  \).
\end{proof}

Central to this proof was finding an \( N  \) normal in \( G  \). But what if we can't? The reason is because there are some groups which do not contain a normal subgroups. 

\begin{definition}[Simple Groups]
   A (finite or infinite) group \( G  \) is called \textbf{simple} if \( | G  |  > 1  \), and the only normal subgroups of \( G  \) are just \( \{ {1}_{G} \}  \) and \( G  \) itself. 
\end{definition}

\begin{eg}[Simple Groups]
       \( Z_p \) where \( p  \) is prime.
\end{eg}
The most important simple group (for us which will be proved later) is called the \textbf{Alternating group}. Consider \( \sigma \in {S}_{3} \) and take \( \sigma = (1 \ 2 \ 3) \) and so we can write this in a couple of ways:
           \[  \sigma = (1 \ 3) (1 \ 2) = (1 \ 2)(1 \ 3)(1 \ 2)(1 \ 3). \]
           Notice that in general for any \( m  \)-cycle in \( {S}_{n} \), we have 
           \begin{align*} 
               &({a}_{1} \ {a}_{2} \ \dots \ {a}_{m}) \\
               &= ({a}_{1} \ {a}_{m}) ({a}_{1} \ {a}_{m-1}) \dots ({a}_{1} \ {a}_{3})({a}_{1} \ {a}_{2}).
           \end{align*}
           That is, every \( m-  \)cycle in \( {S}_{n} \) can be written as a product of \( 2 \)-cycles (which are called transpositions). We know that every permutation in \( {S}_{n} \) can be written as a product of disjoint cycles. Hence, every permutation in \( {S}_{n} \) can be written as a product of transpositions; that is, 
           \[  \langle T  \rangle = {S}_{n} \]
           where \( T = \{  (i \ j) : 1 \leq i < j \leq n  \}  \). As an example, in \( {S}_{4} \), \( T  \) has the elements 
           \[  (1 \ 2), (1 \ 3) ,(1 \ 4), (2 \ 3), (2 \ 4), (3 \ 4). \]

\begin{definition}[Even (Odd) cycles]
    A permutation \( \alpha \in {S}_{n} \) is called \textbf{even} if it can written as a product of an even number of transpositions. Otherwise, \( \alpha  \) is called \textbf{odd}. 
\end{definition}

Recall 
\[  (1 \ 2 \ 3 ) = (1 \ 3)(1 \ 2) \]
and so \( (1 \ 2 \ 3)  \) is even but we also know that 
\[  ( 1 \ 2 \ 3) = (1 \ 2)(1 \ 3)(1 \ 2)(1 \ 3). \]
Is this a well-defined definition? That is, can a permutation be both even and odd? The answer is \textbf{NO}. Permutations in \( {S}_{n} \) are either even or odd but not both. 

\begin{definition}[ ]
   For each \( \sigma \in {S}_{n} \) define 
   \[  \epsilon (\sigma)  = 
   \begin{cases}
       1 &\text{if \( \sigma  \) is even} \\ 
       -1 &\text{if \( \sigma \) is odd}.
   \end{cases} \]
\end{definition}

Notice that \( \epsilon  \) defines a mapping from \( {S}_{n} \) to the multiplicative group \( G = \{  - 1, 1  \}  \) (a subgroup of the quaternion group). Let \( \sigma , \tau \in {S}_{n} \). Assume \( \sigma , \tau  \) are both even or both odd. Then  \( \sigma \tau  \) is an even permutation and so \( \epsilon (\sigma \tau) = 1  \) and \( \epsilon (\sigma) \epsilon (\tau) = 1  \). If one of \( \sigma  \) and \( \tau  \) is even and the other is odd then \( \sigma \tau  \) is odd and so \( \epsilon(\sigma \tau) = -1  \) and \( \epsilon (\sigma) \epsilon(\tau) = -1  \). And so \( \epsilon  \) defines a homomorphism. By the first isomorphism theorem, we have 
\[  {S}_{n} / \text{ker}(\epsilon) \cong \epsilon({S}_{n}) = \{ 1, -1  \}. \]
By Lagrange's theorem, 
\begin{align*}
    \frac{ | {S}_{n} |   }{  | \text{ker}(\epsilon) |  }  = | \{ 1 , -1  \}  | &\implies \frac{ n!  }{ | \text{ker}(\epsilon) |  }  = 2  \\
                                                                               &\implies \frac{ n!  }{  2  }  = |  \text{ker}(\epsilon) |.
\end{align*}
That is, 
\begin{align*}
    \text{ker}(\epsilon) &= \{  \sigma \in {S}_{n} : \epsilon(\sigma)  = 1 \}  \\
                         &= \{  \sigma \in {S}_{n} : \sigma \ \text{is even} \}. 
\end{align*}

\begin{definition}[The Alternating Group is the Kernel of Some homomorphism]
    The alternating group of degree \( n \), denoted by \( {A}_{n} \) is the kernel of homomorphism \( \epsilon  \), but it is more common referred to as the set of all even permutations in \( {S}_{n} \).
\end{definition}


\subsection{Opening Remarks}

\begin{remark}[An Alternative Proof to Problem 1 (a) on Homework]
   If \( \varphi : G \to J  \) where \( \varphi  \) is surjective. If \( J  \) is abelian, then any subgroup of \( G  \) which contains \( \text{ker}(\varphi) \) is normal. 
\end{remark}
\begin{proof}
Note that the first isomorphism theorem says that 
   \[  G / \text{ker}(\varphi)  \cong \varphi(G) \]
   Also, \( \varphi(G) = J  \) since \( \varphi  \) is surjective. So, \( G / \text{ker}(\varphi) \cong J  \) and hence, \( G / \text{ker}(\varphi) \) is abelian. Since \( J  \) is abelian, every subgroup of \( G / \text{ker}(\varphi) \) is normal. Hence, by the 4th iso theorem every subgroup of \( G  \) containing \( \text{ker}(\varphi) \) is also normal.
\end{proof}



\begin{remark}[The Correspondence Theorem]
  Every subgroup of \( G / N  \) is the image of exactly one subgroup of \( G  \) which contains \( N  \). That is, the subgroup of \( G /N  \) is the image of exactly one subgroup of \( G  \) containing \( N  \), through the natural projection. 

  This is to say that if you have a subgroup of \( G / N \), it is of the form  
  \[  \pi(A) = \{  a N : a \in A \} \]
  where \( A \leq G  \) which contains \( N  \). 
\end{remark}

\subsection{Group Actions}

Let \( G  \) be a group and \( A  \) be a nonempty set. A group action is a mapping \( G \times A \to A  \) denoted as 
\[  g \cdot a  \ \  \forall g \in G , \ a \in A. \]
Recall for each \( g \in G  \), there exists a mapping \( \sigma_g : A \to A  \) such that \( g \mapsto g \cdot a  \) which is a permutation or an element of the symmetric group \( {S}_{A} \) on \( A  \). We defined 
\(  \varphi : G \to {S}_{A} \)
defined by \( g \mapsto  \sigma_g \) . We proved that this \( \varphi  \) is a homomorphism. We called it the permutation representation of \( G  \) on \( A  \). Recall further the following important facts:

\begin{enumerate}
    \item[(1)] The kernel of the action is
        \[  \{ g \in G : g \cdot a = a  \ \ \forall a \in A  \} = \text{ker}(\varphi). \]
    \item[(2)] Stabilizer of \( a \in A  \) in \( G  \)
        \[  {G}_{a} = \{  g \in G : g \cdot a = a   \}.  \]
\item[(3)] \textbf{New Fact:} An action is faithful if its kernel is \( \{ {1}_{G} \}  \).
\end{enumerate}

\begin{eg}[Faithful Mapping]
   We know that \( {S}_{n} \) acts on the set \( A = \{  1,2, \dots, n  \}  \) by \( \varphi \cdot i  = \sigma(i) \) for all \( \sigma \in {S}_{n} \) and all \( i \in A  \) where \( \varphi : {S}_{n} \to {S}_{n} \) is the identity map. This is true because \( \sigma(i) = i \) for all \( i \in A  \) and so we have \( \sigma = {1}_{{S}_{n}} \) and so this action is faithful. 
\end{eg}

Recall that \( | {G}_{i} | = (n-1)! \). In general, we can say that \( \text{ker}(\varphi) = \bigcap_{  a \in A  }^{  }  {G}_{a} \). Given a homomorphism \( \psi : G \to {S}_{A} \) where \( A  \) is a nonempty set and \( G  \) is a group, the mapping defined by 
\[  g \cdot a = \psi(g)(a) \] is a group action. Indeed, we see that for all \( {g}_{1}, {g}_{2} \in G  \), we have 
\begin{align*}
    ({g}_{1} {g}_{2}) \cdot a &= \varphi({g}_{1}{g}_{2})(a) \\
                              &= (\varphi({g}_{1}) \circ \varphi({g}_{2})(a)) \\
                              &= \varphi({g}_{1})\Big(\varphi({g}_{2})(a) \Big) \\
                              &= \varphi({g}_{1})({g}_{2} \cdot a ) \\
                              &= {g}_{1} \cdot ({g}_{2} \cdot a )
\end{align*}
and
\[  {1}_{G} \cdot a = \varphi({1}_{G})(a) = {1}_{{S}_{A}}(a) = a.   \]


Notice that 
    \begin{align*}
        \mathcal{O}_a &= \{  b \in A : b \sim a  \}  \\
                      &= \{  b \in A : b = g \cdot a \ \ \text{for some \( g \in G  \)} \} \\
                      &= \{  g \cdot a : g \in G  \}.
    \end{align*}


\begin{prop}[Orbit Stabilizer Theorem]
    Let \( G  \) be a group acting on a nonempty set \( A  \). The relation on \( A  \) defined by \( a \sim b  \) if and only if \( a = g \cdot b  \) for some \( g \in G  \) is an equivalence relation. For each \( a \in A  \), the number of elements in the equivalence class of containing \( a  \) is the index \( | G : {G}_{a} |  \); that is, the number of left-cosets of \( {G}_{a} \) in \( G  \). We denoted the notation for the orbit of \( a  \) as \( \mathcal{O}_a \).    
\end{prop}
\begin{proof}
    To prove the last statement, we need to construct a mapping that is bijective between \( \mathcal{O}_a \) between  

    Let \( \mathcal{O}_a \) be the orbit of \( G  \) containing \( a  \). Let \( b  \in {\mathcal{O}}_{a} \). Then \( b = g \cdot a  \) for some \( g \in G  \). Notice that \( g {G}_{a}  \) is a left-coset of \( {G}_{a} \). Define \( f : {\mathcal{O}}_{a} \to G / {G}_{a} \) by \( b = g \cdot a \longmapsto g {G}_{a} \). Note that we have to be careful here because we might not have a well-defined mapping here. Nevertheless, let us prove that this mapping is onto. Let \( g {G}_{a} \in G / {G}_{a} \). Then \( g \cdot a \in {\mathcal{O}}_{a} \) and \( g \cdot  a  = g {G}_{a} \) and so \(  f \) is surjective. Notice \( f(g \cdot a) = f(h \cdot a) \). Then we have  
    \begin{align*}
        &g {G}_{a} = h {G}_{a} \\
                  &\implies h^{-1} g \in {G}_{a} \\
                  &\implies (h^{-1} g) \cdot a = a \\
                  &\implies h  \cdot (h^{-1} g) \cdot a ) = h \cdot a  \\
                  &\implies (h h^{-1} g ) \cdot a = h \cdot a \\ 
                  &\implies g \cdot a = h \cdot a.
    \end{align*}
    Hence, \( f  \) is injective. But this tells us that \( f  \) is indeed well-defined! Thus, \( f  \) is a bijection from \( {\mathcal{O}}_{a} \) to \( G / {G}_{a} \). Thus, \( | {\mathcal{O}}_{a}  |  = |  G : {G}_{a} |  \). 
\end{proof}

\begin{eg}
   Let \( A = \{ 1,2,3 \}  \) and let \( \sigma \in {S}_{n} \). We know that \( {S}_{n} \) acts on \( A = \{  1,2,\dots, n \}  \). Given \( i \in A  \), find \( {\mathcal{O}}_{i} = \{  \sigma(i) : \sigma \in {S}_{n} \} = A   \). What does the Orbit Stabilizer Theorem say in this example? It tells us that 
   \[  | {\mathcal{O}}_{i} | = | {S}_{n} : {G}_{i}  | \implies n = \frac{ | {S}_{n} |   }{  | {G}_{i} |   } = \frac{ n!  }{ | {G}_{i} |  } \implies | {G}_{i} | = \frac{ n!  }{  n  }  = (n-1)!.  \]
\end{eg}

